\section{Fundamental Particles}
    \subsection{Quarks (\texorpdfstring{$\quark$}.)}
        All quarks have non-zero mass and charge.
        They all have a baryon number \\ of $B=+\frac{1}{3}$.
        All antiquarks ($\antiquark$) have the opposite charge of their regular counterparts, and a baryon number of $B=-\frac{1}{3}$.
        The quarks are split into generations, $\upquark$ and $\downquark$ are generation I, $\strangequark$ and $\charmquark$ are in generation II and $\topquark$ and $\bottomquark$ are in generation III. \\

        \begin{center}
        \begin{tabular}{l l c c c}
            \textbf{Quark}            & \textbf{Antiquark}                & \textbf{Mass/}MeV & \textbf{Charge/}$e$ & \textbf{Strangeness} \\[0.4em]
            Up ($\upquark$)           & Antiup ($\antiupquark$)           &                   & $+\frac{2}{3}$ & 0 \\[0.4em]
            Down ($\downquark$)       & Antidown ($\antidownquark$)       &                   & $-\frac{1}{3}$ & 0 \\[0.4em]
            Strange ($\strangequark$) & Antistrange ($\antistrangequark$) &                   & $+\frac{2}{3}$ & -1 \\[0.4em]
            Charm ($\charmquark$)     & Anticharm ($\anticharmquark$)     &                   & $-\frac{1}{3}$ & 0 \\[0.4em]
            Top ($\topquark$)         & Antitop ($\antitopquark$)         &                   & $+\frac{2}{3}$ & 0 \\[0.4em]
            Bottom ($\bottomquark$)   & Antibottom ($\antibottomquark$)   &                   & $-\frac{1}{3}$ & 0
        \end{tabular}
        \end{center}

    \subsection{Leptons (\texorpdfstring{$\lepton$}.)}
        All leptons have a non-zero mass.
        The charged leptons ($\electron[-]$, $\muon[-]$, $\taulepton[-]$) all have a charge of -1; the neutrinos ($\neutrino[e]$, $\neutrino[\mu]$, $\neutrino[\tau]$) are all neutral.
        All leptons have a lepton number of $\leptonnumber[]=+1$, or more specifically, $\electron[-]$ and $\neutrino[e]$ have an electron number of $\leptonnumber[e]=+1$; $\muon[-]$ and $\neutrino[\mu]$ have a muon number of $\leptonnumber[\mu]=+1$ and $\taulepton[-]$ and $\neutrino[\tau]$ have a tau number of $L_\tau=+1$.
        All antileptons ($\antilepton$) have the opposite charge of their regular counterparts (this is obviously unchanged for neutrinos, since they are neutral), and their lepton numbers are $\leptonnumber=-1$. \\

        \begin{center}
        \begin{tabular}{l l c c}
            \textbf{Lepton}                    & \textbf{Antilepton}                        & \textbf{Mass/}MeV & \textbf{Charge/}$e$ \\
            Electron ($\electron[-]$)          & Positron ($\electron[+]$)                  &                   & -1                  \\
            Electron Neutrino ($\neutrino[e]$) & Electron Antineutrino ($\antineutrino[e]$) &                   & 0                   \\
            Muon ($\muon[-]$)                  & Antimuon ($\muon[+]$)                      &                   & -1                  \\
            Muon Neutrino ($\neutrino[\mu]$)   & Muon Antineutrino ($\antineutrino[\mu]$)   &                   & 0                   \\
            Tau ($\taulepton[-]$)              & Antitau ($\taulepton[+]$)                  &                   & -1                  \\
            Tau Neutrino ($\neutrino[\tau]$)   & Tau Antineutrino ($\antineutrino[\tau]$)   &                   & 0
        \end{tabular}
        \end{center}

    \subsection{Bosons}
        All bosons except the Higgs boson ($\higgsboson$) mediate one of the fundamental forces:
        \begin{itemize}
            \item Photons ($\photon$) mediate the electromagnetic force
            \item Gluons ($\gluon$) mediate the strong nuclear force
            \item $\wboson[]$ and $\zboson$ bosons ($\wboson[\pm]$ and $\zboson$) mediate the weak nuclear force
            \item Gravitons ($\graviton$) mediate gravity
        \end{itemize}
        
        Gravitons are hypothetical. \\

        \begin{center}
        \begin{tabular}{l l c c}
            \textbf{Boson}           & \textbf{Antiboson}       & \textbf{Mass/}MeV & \textbf{Charge/}$e$ \\
            Photon ($\photon$)       & Photon ($\photon$)       &                   & +1                  \\
            Gluon ($\gluon$)         & Gluon ($\gluon$)         &                   & 0                   \\
            W Boson ($\wboson[\pm]$) & W Boson ($\wboson[\mp]$) &                   & $\pm1$              \\
            Z Boson ($\zboson$)      & Z Boson ($\zboson$)      &                   & 0                   \\
            Higgs ($\higgsboson$)    & Higgs ($\higgsboson$)    &                   & 0                   \\
            Graviton ($\graviton$)   & Graviton ($\graviton$)   &                   & 0
        \end{tabular}
        \end{center}

    \newpage

\section{Mesons}
    There are 2 types of mesons required for A level physics: pions ($\pion$) and kaons ($\kaon$).
    Kaons contain strange quarks ($\strangequark$) or antistrange quarks ($\antistrangequark$), so they have a strangeness of $\pm1$. \\

    \begin{center}
    \begin{tabular}{l l l}
        \textbf{Meson}                    & \textbf{Quark Flavour}                                    & \textbf{Antimeson}                \\
        Positive Pion ($\pion[+]$)        & $\mesonquark[u][d]$                                       & Negative Pion ($\pion[-]$)        \\
        Negative Pion ($\pion[-]$)        & $\mesonquark[d][u]$                                       & Positive Pion ($\pion[+]$)        \\
        Neutral Pion ($\pion[0]$)         & $\frac{1}{\sqrt{2}}(\mesonquark[u][u]-\mesonquark[d][d])$ & Neutral Pion ($\pion[0]$)         \\
        Positive Kaon ($\kaon[+]$)        & $\mesonquark[u][s]$                                                & Negative Kaon ($\kaon[-]$)        \\
        Negative Kaon ($\kaon[-]$)        & $\mesonquark[s][d]$                               & Positive Kaon ($\kaon[+]$)        \\
        Neutral Kaon ($\kaon[0]$)         & $\mesonquark[d][s]$                             & Neutral Antikaon ($\antikaon[0]$) \\
        Neutral Antikaon ($\antikaon[0]$) & $\mesonquark[s][u]$                               & Neutral Kaon ($\kaon[0]$)
    \end{tabular} \\

    \begin{tabular}{l c c c}
        \textbf{Meson} & \textbf{Mass/}MeV & \textbf{Charge/}$e$ & \textbf{Strangeness} \\
        $\pion[+]$        &                   & +1                  & 0                 \\
        $\pion[-]$        &                   & -1                  & 0                 \\
        $\pion[0]$        &                   & 0                   & 0                 \\
        $\kaon[+]$        &                   & +1                  & -1                \\
        $\kaon[-]$        &                   & -1                  & +1                \\
        $\kaon[0]$        &                   & 0                   & -1                \\
        $\antikaon[0]$    &                   & 0                   & +1
    \end{tabular}
    \end{center}
    \newpage

\section{Baryons}
    All baryons are made up of three quarks, meaning that they have a baryon number of $B=+1$.
    Antibaryons thus have a baryon number of $B=-1$, since they are made up of three antiquarks.
    Only nucleons ($\nucleon$) (\ref{nucleons}) and sigma baryons ($\sigmabaryon$) (\ref{sigma}) are shown below, but there are also delta, lambda, xi, and omega baryons ($\deltabaryon$, $\lambdabaryon$, $\varXi$, $\omegabaryon$).
    There are possibly also `exotic baryons' with more quarks and antiquarks; these are hypothetical and not also not shown below. \\

    \subsection{Nucleons (\texorpdfstring{$\nucleon$}.)} \label{nucleons}
        Nucleons are the most common baryons, since they are the the components of atoms. \\

        \begin{center}
        \begin{tabular}{l l l}
            \textbf{Baryon}                    & \textbf{Quark Flavour} & \textbf{Antibaryon}                        \\
            Proton ($\proton$/$\nucleon[+]$)   & $\baryonquarks[uud]$   & Antiproton ($\antiproton$/$\nucleon[-]$)   \\
            Neutron ($\neutron$/$\nucleon[0]$) & $\baryonquarks[udd]$   & Antineutron ($\antineutron$/$\nucleon[0]$)
        \end{tabular}
        \end{center}

        \begin{center}
        \begin{tabular}{l c c c}
            \textbf{Proton} & \textbf{Mass/}MeV & \textbf{Charge/}$e$ & \textbf{Strangeness} \\
            $\proton$       &                   & +1                  & 0 \\
            $\neutron$      &                   & 0                   & 0 \\
        \end{tabular}
        \end{center}

    \subsection{Sigma (\texorpdfstring{$\sigmabaryon$}.)} \label{sigma}
        Sigma baryons ($\sigmabaryon$) are made up of two generation I quarks ($\upquark$ or $\downquark$), and the third quark is either generation II or III ($\strangequark$, $\charmquark$, $\topquark$, $\bottomquark$).
        Only the strange sigmas are shown below. \\

        \begin{center}
        \begin{tabular}{l l l}
            \textbf{Baryon}                    & \textbf{Quark Flavour} & \textbf{Antibaryon} \\
            Positive Sigma ($\sigmabaryon[+]$) & $\baryonquarks[uus]$   & Negative Antisigma ($\antisigmabaryon[-]$) \\
            Negative Sigma ($\sigmabaryon[-]$) & $\baryonquarks[uds]$   & Positive Antisigma ($\antisigmabaryon[+]$) \\
            Neutral Sigma ($\sigmabaryon[0]$)  & $\baryonquarks[dds]$   & Neutral Antisigma ($\antisigmabaryon[0]$)
        \end{tabular}
        \end{center}

        \begin{center}
        \begin{tabular}{l c c c}
            \textbf{Meson}        & \textbf{Mass/}MeV & \textbf{Charge/}$e$ & \textbf{Strangeness} \\
            $\sigmabaryon[+]$     &                   & +1                  & -1 \\
            $\sigmabaryon[-]$     &                   & -1                  & -1 \\
            $\sigmabaryon[0]$     &                   & 0                   & -1 \\
        \end{tabular}
        \end{center}